\documentclass[11pt,a4paper]{article}

\usepackage[margin=1in]{geometry}
\usepackage{amsmath,amssymb}
\usepackage{booktabs}
\usepackage{listings}
\usepackage{xcolor}
\usepackage{hyperref}
\usepackage{enumitem}
\usepackage[expansion=false]{microtype}
\usepackage{titlesec}
\usepackage{verbatim}

\hypersetup{
    colorlinks=true,
    linkcolor=blue!50!black,
    urlcolor=blue!50!black,
    citecolor=blue!50!black,
}

\definecolor{codebg}{rgb}{0.97,0.97,0.97}
\definecolor{codekw}{rgb}{0.20,0.20,0.60}
\definecolor{codecm}{rgb}{0.40,0.55,0.40}
\definecolor{codestr}{rgb}{0.55,0.30,0.30}

\lstset{
    backgroundcolor=\color{codebg},
    basicstyle=\ttfamily\small,
    keywordstyle=\color{codekw}\bfseries,
    commentstyle=\color{codecm}\itshape,
    stringstyle=\color{codestr},
    showstringspaces=false,
    breaklines=true,
    frame=single,
    framesep=4pt,
    rulecolor=\color{black!20},
    aboveskip=8pt,
    belowskip=8pt,
    columns=fullflexible,
    keepspaces=true,
}

\lstdefinestyle{cstyle}{language=C}
\lstdefinestyle{plainstyle}{language=}

\titleformat{\section}{\Large\bfseries}{\thesection}{1em}{}
\titleformat{\subsection}{\large\bfseries}{\thesubsection}{1em}{}

\title{\texttt{simdq} Plan 1 --- The Science of Bit-Hashed Retrieval}

\author{DocUVerse \texttt{simdq} engine project}
\date{2026-06-17}

\begin{document}
\maketitle

\begin{center}
\small
\textbf{Project:} DocUVerse \texttt{simdq} engine \\
\textbf{Scope of this paper:} Plan 1 --- the C/SIMD kernel layer (no Python, no engine wrapper) \\
\textbf{Branch:} \texttt{v0.1.2} \\
\textbf{Companion documents:} design spec, plan \& status, bench data
\end{center}

\vspace{1em}

\section{Abstract}

Plan 1 of the \texttt{simdq} retrieval engine ports a 768-bit symmetric
Hamming top-1 kernel into the DocUVerse tree and generalises it along
two axes:

\begin{enumerate}
    \item \textbf{Top-K} --- a bounded max-heap is threaded through the
    SIMD inner loop without breaking its straight-line shape, supporting
    $K \le 256$ on the stack with a single thresholded compare per
    \texttt{LANES}-block.
    \item \textbf{Asymmetric $b$-bit scans} --- full-precision query
    vectors scored against $b \in \{1, 2, 4\}$ quantized database codes,
    holding the query's bits intact so the score estimator's variance
    scales with the \emph{square} of the level spacing rather than the
    floor of a sign-only Hamming match.
\end{enumerate}

This document explains the \emph{information-theoretic} and
\emph{architectural} reasons those design choices land where they do
--- i.e.\ the science behind the bit hashing --- and ties each piece
back to the kernel code that ships in
\path{docuverse/engines/retrieval/simdq/_native/}.

\section{Why bit-hashed retrieval works}

A trained dense embedding model produces vectors $x \in \mathbb{R}^D$
whose cosine similarity is meaningful for retrieval. For modern
encoders (e.g.\ \texttt{granite-embedding-311m-multilingual-r2},
$D = 768$), the \textbf{sign pattern} of $x$ already carries most of
that information.

\paragraph{The Charikar SimHash bound.} Draw a single isotropic random
hyperplane $h \in \mathbb{R}^D$ (i.e.\ $h \sim \mathcal{N}(0, I_D)$).
For unit vectors $x, y$:
\begin{equation}
\Pr_{h}\bigl[\,\mathrm{sign}(\langle h, x\rangle) = \mathrm{sign}(\langle h, y\rangle)\,\bigr]
\;=\; 1 - \frac{\theta(x,y)}{\pi},
\end{equation}
where $\theta(x, y) = \arccos(\langle x, y\rangle)$. This is a
\emph{per-hyperplane} statement: the probability that one random
hyperplane separates $x$ and $y$ is $\theta/\pi$, equivalently the
probability they agree is $1 - \theta/\pi$. (It is \emph{not} the
joint probability that all $D$ hyperplanes agree --- that would be
$(1 - \theta/\pi)^D$, which is vanishingly small at $D = 768$.)

Now draw $D$ independent hyperplanes $h_1, \dots, h_D$ and hash $x,
y$ into $D$-bit codes by taking the signs of the $D$ projections.
Let $A_w = \mathbf{1}[\mathrm{sign}(\langle h_w, x\rangle) =
\mathrm{sign}(\langle h_w, y\rangle)]$ be the indicator that bit $w$
agrees. By Eq.~(1), $\mathbb{E}[A_w] = 1 - \theta/\pi$ for each $w$,
so by linearity of expectation:
\begin{itemize}
  \item the \textbf{expected fraction of agreeing bits} is
  $\mathbb{E}\bigl[\tfrac{1}{D}\sum_w A_w\bigr] = 1 - \theta/\pi$
  --- a deterministic, monotone function of cosine;
  \item equivalently --- and this is the kernel-relevant form --- the
  \textbf{Hamming distance} $\sum_w (1 - A_w)$ is an unbiased
  estimator of $(D/\pi)\cdot\theta(x,y)$.
\end{itemize}

For $D = 768$, the standard deviation of that estimator at a target
similarity $\cos\theta \approx 0.9$ is roughly
\begin{equation}
\sigma_H \;\approx\; \sqrt{D \cdot p(1-p)} \;\approx\; \sqrt{768 \cdot 0.146 \cdot 0.854} \;\approx\; 10\ \text{bits},
\end{equation}
small relative to the dynamic range of plausible Hamming distances
($\approx 0$--$384$), and small relative to the typical per-document
cosine gap between true relevance and a near-miss. That is the
empirical reason 1-bit codes work at all for text retrieval, and it is
why the existing \texttt{hb5} driver achieves competitive recall on
BEIR without any rescore tier.

\paragraph{The full-precision encoder is its own random projection.}
When the encoder is trained well, the dense embedding lives on a
low-curvature manifold inside $\mathbb{R}^{768}$, and the directions
sampled by the trained model already act like a useful random
projection of the underlying semantic space. Plan 1 exploits this
directly: it does \textbf{no extra projection} for the symmetric
kernel --- $\mathrm{sign}(x)$ is the code, full stop. (The spec keeps
random orthogonal $W$ in scope for later $(b, d)$ recipes; Plan 1
just doesn't need it.)

\section{The symmetric kernel: XOR-popcount on SoA codes}

The Hamming distance between two $\texttt{WORDS} = 12$ $\times$ 64-bit
codes (768 bits) is
\begin{equation}
H(q, x) \;=\; \sum_{w=0}^{\texttt{WORDS}-1} \mathrm{popcnt}(q_w \oplus x_w).
\end{equation}
This is the inner loop of \texttt{scan\_soa512} in
\path{simdq_kernels_hamming.h:52-77} and of \texttt{scan\_shard} in
\path{simdq_kernels_hamming.h:87-123} (AVX-512 path) /
\path{simdq_kernels_hamming.h:152-195} (AVX2 path).

\subsection{Why SoA, not AoS}

Two layouts are admissible (\path{simdq_common.h:7-8}):

\begin{lstlisting}[style=plainstyle]
AoS:  db [ i*WORDS + w ]   -- code i contiguous, words sequential
SoA:  dbT[ w*n + i ]       -- word w contiguous across all codes
\end{lstlisting}

For a vertical SIMD scan, the inner loop wants to load \textbf{the same
word $w$ of \texttt{LANES} consecutive codes} in a single 256/512-bit
load. That's $\texttt{dbT}[w \cdot n + i \,..\, w \cdot n + i +
\texttt{LANES} - 1]$ --- a unit-stride load --- under SoA, but a
strided gather under AoS ($\texttt{WORDS} \cdot 8 = 96$ byte stride).
On x86-64 the gather is bandwidth-bound at roughly one element per
cycle, the contiguous load is bound at 32--64 bytes per cycle. The
factor is exactly the SIMD width.

The cost of SoA is paid once at ingest (a transpose). The cost of AoS
is paid every query, forever.

\subsection{The two SIMD paths}

\paragraph{AVX-512 + VPOPCNTDQ (Zen 4/5, Ice Lake+).}
\path{simdq_kernels_hamming.h:54-69}:
\begin{lstlisting}[style=cstyle]
acc += _mm512_popcnt_epi64(_mm512_xor_si512(d, broadcast(q[w])));
\end{lstlisting}
8 codes per iteration, 12 word-loads per block, \textbf{96 popcnt-bits
per vpopcnt instruction}. The accumulator stays in a \texttt{\_\_m512i}
of 8 $\times$ \texttt{u64} lane counters; the running best per lane is
updated with \texttt{\_mm512\_cmplt\_epi64\_mask} +
\texttt{\_mm512\_mask\_mov\_epi64}, eliminating the last branch from
the hot loop.

\paragraph{AVX2 nibble-LUT (Haswell..Zen 3).}
\path{simdq_kernels_hamming.h:135-140}:
\begin{equation}
\mathrm{popcnt\_byte}(b) \;=\; \mathrm{shuffle}(\mathrm{LUT}, b \,\&\, \mathtt{0x0F}) + \mathrm{shuffle}(\mathrm{LUT}, (b \gg 4) \,\&\, \mathtt{0x0F}).
\end{equation}
\texttt{vpshufb} does a 16-entry table lookup per byte lane in one
cycle; the two halves of each byte (two nibbles) are looked up
independently and summed. The byte counters then accumulate across all
12 words --- the bound is $12 \cdot 8 = 96 < 255$, so an 8-bit
accumulator cannot overflow. One \texttt{vpsadbw} per 4-code block
reduces 32 bytes of partial counts to 4 $\times$ \texttt{u64}
distances.

This nibble-LUT trick is the textbook AVX2 popcount; what matters here
is that \textbf{the byte-overflow argument is exact}, not heuristic ---
the kernel exploits the fact that no 8-bit lane can ever exceed 96 to
skip a per-word reduction.

\subsection{Hamming top-K is essentially free}
\label{sec:topk-free}

\path{simdq_kernels_hamming_topk.h:46-78} adds top-K to the AVX-512
kernel in eleven lines: after each \texttt{LANES}-block, the per-lane
distances are spilled to a small scratch array, compared against a
single \texttt{int64\_t threshold = heap.keys[0]} (the largest of the
$K$ kept), and offered to \texttt{simdq\_topk\_offer} only if they
pass. For random data and $K = 100$ against $n = 50\,\mathrm{M}$, the
threshold settles fast and the heap-offer slow path almost never runs
--- the bench data show \texttt{bench\_hamming} at \textbf{740 M
cmp/s} single-query against \texttt{bench\_hamming\_top1} at
\textbf{288 M cmp/s per query} in the \texttt{NQ=8} batch, so
single-query top-K with $K = 100$ is \emph{faster} than per-query
top-1 in the \texttt{NQ=8} batch. The win is register-pressure: top-K
keeps fewer in-flight accumulators than the 8-query batched top-1.

\subsection{The query-batching trick (carried over from \texttt{hb5})}

Memory bandwidth, not compute, sets the ceiling for an XOR-popcount
loop on a Zen 4-class chip. The cure, due to the existing \texttt{hb5}
work, is to \textbf{load each database word once and reuse it for
\texttt{NQ} queries}:

\begin{lstlisting}[style=cstyle]
for (w = 0; w < WORDS; w++) {
    d = load(dbT + w*n + i);            // ONE load
    for (j = 0; j < NQ; j++)             // NQ XOR+popcounts
        acc[j] += popcnt(d XOR qs[j][w]);
}
\end{lstlisting}

(\path{simdq_kernels_hamming.h:99-115}). The arithmetic stays the
same; the memory traffic is divided by \texttt{NQ}. At
$\texttt{NQ} = 8$ and $n = 50\,\mathrm{M}$, the AVX2 nibble-LUT path
hits \textbf{2308 M cmp/s @ 27.7 GB/s} --- within a few percent of the
DRAM streaming ceiling on the dev host.

\section{The asymmetric kernels: information beats bits}

\subsection{Why asymmetric}

If we quantise \emph{both} the query and the database to $b$ bits, the
score we recover is a noisy estimate of an inner product between two
quantised vectors. The variance of that estimate is set by the
\emph{square} of the quantization step on \textbf{both sides} --- and
for $b = 1$ the database side already loses an enormous amount of
magnitude information (a unit-norm $x$ collapses to $\pm 1$ per dim).

If instead we keep the \textbf{query in fp32} and quantise only the
database, the score
\begin{equation}
\hat{s}_i \;=\; \sum_{w} q'_w \cdot v_{i,w}, \qquad v_{i,w} \in V_b
\end{equation}
is an unbiased estimate of $q' \cdot y_i$ with variance
\begin{equation}
\mathrm{Var}[\hat{s}_i] \;=\; \sum_{w} (q'_w)^2 \cdot \mathrm{Var}[v_{i,w} - y_{i,w}]
\;\le\; \|q'\|^2 \cdot \frac{\Delta_b^2}{12},
\end{equation}
(per-dim quantization error variance, with $\Delta_b$ the level
spacing). For $V_2 = \{-3, -1, +1, +3\}$ after the per-vector scale
$s_i = \|y_i\|/\sqrt{d}$ is applied, $\Delta_b = 2$; for $V_4$,
$\Delta_b = 2$ as well over a 16-level grid. The \textbf{query being in
float means its directional information is fully preserved} --- only
the database codebook's projection error contributes noise.

This is the core ASH (Asymmetric Scalar Hashing) observation, and it
is the reason Plan 1 ships three asymmetric kernels rather than just
the symmetric Hamming one.

\subsection{Per-vector scale}

\path{simdq_pack.h:32-44}:
\begin{equation}
s_i \;=\; \frac{\|y_i\|}{\sqrt{d}}
\end{equation}
stored as one fp16 (or fp32 in Plan 1 for now) per database vector.
This is the unique scalar that makes the quantization grid's mean
squared per-dim residual minimal \emph{under the assumption that $y_i$
has roughly isotropic dimensions}. After dividing by $s_i$, each
$y_{i,w}$ has unit MSE, so the level grid $V_b = \{2c - (2^b - 1) \mid
c = 0\dots 2^b - 1\}$ catches roughly the right fraction of the
distribution at every boundary --- i.e.\ for $b = 2$ the levels
$\{-3, -1, +1, +3\}$ straddle the fitted standard normal at the
natural quartile boundaries.

The cost of this scale is \textbf{2 bytes per code} at fp16 --- under
0.5\% of a $b = 4$, $d = 768$ code, and 2\% of a $b = 1$ code. Recall
gain in published ASH-style work is a few NDCG points, so the trade is
unambiguously worth it.

\subsection{$V_b$ is the natural symmetric grid}

The grid $V_b = \{2c - (2^b - 1) \mid c = 0\dots 2^b - 1\}$ has three
properties that make it the right choice for an asymmetric scan:

\begin{enumerate}
    \item \textbf{Symmetric around 0.} No bias on a zero-mean coordinate
    distribution.
    \item \textbf{Integer-valued.} \texttt{vf[l] = (float)levels[code]}
    is exact, no rounding noise injected on top of the quantization
    step.
    \item \textbf{Maps to $c \in \{0, \dots, 2^b - 1\}$ with one
    shift+mask.} The encoded ``code'' is a $b$-bit unsigned that the
    SIMD path can extract with \texttt{(packed >> (l*b)) \& ((1 << b) -
    1)} --- no per-lane sign extension on the AVX2 path, which would
    otherwise cost a separate shuffle.
\end{enumerate}

The alternative --- a Lloyd--Max grid fit to the empirical coordinate
distribution --- would gain at most a fraction of a bit of effective
resolution and would forfeit (3). Plan 1 doesn't bother.

\subsection{The packed SoA layout, generalised}

For asymmetric scans the SoA stride is \textbf{$d \times \lceil N
\cdot b / 8 \rceil$ bytes}, with each dim row holding $N$ codes packed
$b$ bits wide (\path{simdq_pack.h:1-19}):

\begin{lstlisting}[style=plainstyle]
b = 1 :  byte = i >> 3,   bit   = i & 7
b = 2 :  byte = i >> 2,   shift = 2 * (i & 3)
b = 4 :  byte = i >> 1,   shift = 4 * (i & 1)
\end{lstlisting}

This is dim-major because the inner loop, exactly as for the Hamming
kernel, broadcasts a single $q'_w$ and FMAs it across \texttt{LANES}
consecutive codes' decoded levels. The \texttt{LANES} is set by the
SIMD register width: $\texttt{ASYM\_LANES} = 16$ for AVX-512
(\texttt{\_\_m512}), $\texttt{ASYM\_LANES} = 8$ for AVX2
(\texttt{\_\_m256}).

\subsection{Inner loop, $b=2$ and $b=4$ (the fast paths)}

\path{simdq_kernels_asym_b2.h:35-62} (AVX-512F) and lines 99--121
(AVX2-FMA) implement the $b = 2$ scan;
\path{simdq_kernels_asym_b4.h} is structurally identical at $b = 4$.
Sketch:

\begin{lstlisting}[style=plainstyle]
for each block of LANES codes:
    acc = 0
    for w = 0 .. d-1:
        packed = *(uintN_t*)(codes + w*row_bytes + (i0 >> log2(8/b)))
        for l = 0 .. LANES-1:
            vf[l] = levels[ (packed >> (l*b)) & ((1<<b)-1) ]
        acc = fmadd( broadcast(q[w]), load(vf), acc )
    spill acc to LANES floats; threshold-test against heap;
        offer survivors.
\end{lstlisting}

The 16-byte LUT for $b = 2$ ($\{-3, -1, +1, +3\}$ repeated to fill
\texttt{vpshufb}) and the 16-byte LUT for $b = 4$ ($\{2c - 15 \mid c =
0\dots 15\}$) are immediates the compiler hoists out of the loop on
AVX-512 hosts.

The benchmark numbers (single-threaded, 50~M codes):

\begin{center}
\begin{tabular}{cccc}
\toprule
$b$ & bytes/code & cmp/s & GB/s \\
\midrule
2 & 192 & 2.19 M & 0.4 \\
4 & 384 & 2.36 M & 0.9 \\
\bottomrule
\end{tabular}
\end{center}

are bound by the FMA throughput of the AVX2 path on the dev host, not
by memory. On an AVX-512F-capable chip the \texttt{\_\_m512}
accumulator doubles the per-cycle FMA work and the per-iteration
unpack drops from a scalar 8-step expansion to a single masked-blend
or shuffle --- we expect roughly \textbf{2--3$\times$ the AVX2
throughput} on Zen 4 / Sapphire Rapids, which is the next measurement
Plan 2 must take.

\subsection{The $b=1$ asymmetric outlier}
\label{sec:b1-outlier}

\path{simdq_kernels_asym_b1.h} is the most informative kernel in
Plan 1 because the \textbf{AVX-512F path is provably good} while the
\textbf{AVX2 path is provably suboptimal}:

\begin{itemize}
    \item \textbf{AVX-512F (lines 46--72).} 16 packed bits $\to$
    \texttt{\_\_mmask16} $\to$ \texttt{\_mm512\_mask\_blend\_ps(m, -1,
    +1)} in \textbf{one instruction}. Then a single FMA against the
    broadcast query. This is dense, branchless, and bandwidth-bound
    --- exactly the design point.

    \item \textbf{AVX2-FMA (lines 99--130).} The mask register doesn't
    exist below AVX-512, so the kernel currently expands 8 bits to 8
    floats with a per-iteration scalar loop:
\begin{lstlisting}[style=cstyle]
for (l = 0; l < 8; l++)
    vf[l] = (bits & (1u << l)) ? 1.0f : -1.0f;
\end{lstlisting}
    The compiler does \emph{not} vectorize this in a useful way. The
    observed ratio at $n = 50\,\mathrm{M}$ (\texttt{bench\_asym\_b1})
    is \textbf{0.55 M cmp/s}, $\sim 6\times$ slower than $b = 2$
    despite reading 2$\times$ \emph{less} memory.
\end{itemize}

The AVX2 fix is a known three-instruction trick:
\texttt{\_mm256\_set1\_epi8(bits)} $\to$ \texttt{vpshufb} against an
8$\times$\texttt{u8} spread table $\to$ masked blend to $\pm 1$
floats. This is queued for Plan 2. Plan 1's value here is that
\textbf{the bottleneck is now measured and named} --- it would have
been hidden under ``the kernel is just slow'' without a microbench
targeting exactly this path.

\subsection{The bandwidth budget}

For asymmetric $b = 2$, $d = 768$, the per-query work is
\begin{align}
\text{loads}      &: d \cdot \lceil N \cdot b / 8 \rceil = 768 \cdot N/4 = 192\,N \text{ bytes} \\
\text{arithmetic} &: d \cdot \texttt{LANES} \text{ FMAs} / \texttt{LANES} \text{ codes} \approx 768 \text{ FMAs per code}
\end{align}

On Zen 4 (single-threaded, $\sim 64$ GB/s sustained DRAM, 2 FMA/cycle
at $\sim 3.5$ GHz on \texttt{\_\_m512}):
\begin{align}
\text{bandwidth ceiling} &: \frac{64 \text{ GB/s}}{192 \text{ B/code}} = 333 \text{ M cmp/s} \\
\text{compute ceiling}   &: \frac{16 \text{ lanes} \cdot 2 \text{ FMA/cyc} \cdot 3.5 \text{ GHz}}{768 \text{ FMA/code}} = 146 \text{ M cmp/s}
\end{align}

so on Zen 4 the \textbf{$b = 2$ kernel is compute-bound}, not
memory-bound, and the ceiling is reached when both FMA pipes stay
busy. The AVX-512 path (16 lanes vs 8) doubles compute headroom; the
AVX2 path is half. This matches the empirical AVX2 numbers
($\sim 2.2$ M cmp/s $\times$ 32 threads $\approx 70$ M cmp/s, within a
factor of 2 of the AVX2-compute ceiling once the threaded driver
lands in Plan 2).

For $b = 1$, the load drops to $96 N$ bytes per query.
\textbf{Memory is no longer the bottleneck}; the FMA pipe count is.
This is why Hamming + popcount is faster than asymmetric $b = 1$
despite reading the same bytes --- popcount uses a different
instruction port from the FMA pipes, so the symmetric kernel's
effective issue rate is higher.

\section{Top-K extraction in the SIMD inner loop}

\path{simdq_topk.h} is a 100-line bounded max-heap with two design
constraints:

\begin{enumerate}
    \item \textbf{Caller-allocated.} The \texttt{keys} and \texttt{idxs}
    arrays are passed in; the heap struct itself is 24 bytes and lives
    on the stack of the scan kernel. No allocation in the hot path.
    \item \textbf{Branchless thresholding.}
    \texttt{simdq\_topk\_threshold(h)} returns \texttt{INT64\_MAX}
    while the heap is filling and \texttt{h->keys[0]} once full. The
    inner-loop test is therefore always a single integer compare; the
    heap-touching code (\texttt{simdq\_topk\_offer}) only runs on
    survivors.
\end{enumerate}

For $K \le 256$ random codes per shard against $n \approx 1\,\mathrm{M}$,
the survival rate per \texttt{LANES}-block converges to roughly $K /
n$ after the first few percent of the scan, so the amortised heap
cost is \textbf{well under one offer per block}. The empirical
evidence is in \S\ref{sec:topk-free}: top-K at $K = 100$ runs at the
same per-query throughput as top-1.

\subsection{Negated fixed-point keys for asymmetric scores}

Asymmetric scans want the \textbf{$K$ largest} scores; the heap
supports \textbf{$K$ smallest} keys. The kernels resolve this with a
fixed-point negation (see e.g.\ \path{simdq_kernels_asym_b2.h:56}):
\begin{equation}
\text{key} \;=\; -\,(\texttt{int64\_t})(\,\text{score} \cdot 2^{20}\,)
\end{equation}
The $2^{20}$ factor preserves $\sim 6$ decimal digits of the float
score in an \texttt{int64} --- more than enough to break ties
correctly between any two non-equal asymmetric inner products on a
$d = 768$ corpus. The choice of integer keys (rather than float keys)
keeps the heap comparison a single \texttt{cmp + jl} instruction,
which matters because the heap is touched inside the inner loop.

\subsection{OpenMP merge for the threaded Hamming driver}

\path{simdq_kernels_hamming_topk.h:143-168} runs $T$ shards in
parallel, each with its own per-thread heap, and merges into a global
heap under \texttt{omp critical}. For $T = 32$ and $K = 100$, the
merge work is $T \cdot K \cdot \log K \approx 32 \cdot 100 \cdot 7
\approx 22\,000$ operations --- a rounding error next to a 50 M-code
scan. Plan 1 ships this for the Hamming family only; the asymmetric
scans inherit the same pattern in Plan 2.

The single subtlety is \textbf{partial-fill safety}: a thread whose
shard is shorter than $K$ will not fill its local heap to $K$
entries, so the merge code initialises \texttt{ld[r] = INT64\_MAX,
li[r] = -1} and skips sentinel rows during merge (line 159,
\texttt{if (li[r] >= 0 \&\& \dots)}). This is the kind of bug that
doesn't show up at $n \gg T \cdot K$ and bites exactly once --- at
the boundary where some user runs $n = 50$ for a unit test.

\section{Test coverage and what it actually checks}

Plan 1's 12 ctest targets aren't ticking a checkbox; they are designed
to falsify a specific class of kernel bug each:

\begin{center}
\begin{tabular}{ll}
\toprule
\textbf{Target} & \textbf{Falsifies} \\
\midrule
\texttt{kernels\_hamming\_*} & XOR-popcount mismatch vs scalar reference \\
\texttt{kernels\_hamming\_topk\_*} & top-K boundary (K=1, planted, threaded shard boundary) \\
\texttt{kernels\_asym\_b\{1,2,4\}\_*} & level grid mis-mapping, per-vector scale, planted-best \\
\texttt{pack} & round-trip pack$\to$unpack for all $b$, including b=4 saturation clamp \\
\texttt{topk} & heap invariants under random offers; ascending extraction \\
\bottomrule
\end{tabular}
\end{center}

The two SIMD paths are exercised independently. On an AVX-512F host,
\texttt{*\_native} compiles with \texttt{-march=native} (picks
AVX-512), \texttt{*\_avx2} forces \texttt{-mavx2 -mpopcnt} (or
\texttt{-mavx2 -mfma} for asymmetric). On an AVX2-only host,
\texttt{*\_native} \emph{is} \texttt{*\_avx2} and the AVX-512 paths
are not exercised --- which is why \S\ref{sec:results} below flags
re-running the sweep on an AVX-512F box as a Plan-2 prerequisite.

The planted-best tests are the most targeted: they insert the known
best at index 0, at index $N - 1$, in the $n \bmod \texttt{LANES}$
tail, and across an OpenMP shard boundary. This is the test that
catches the class of bug where the SIMD body is correct but the tail
loop or the cross-thread merge silently drops the winning index.

\section{Empirical results and what they mean}
\label{sec:results}

Headline numbers from \texttt{\_native/BENCHMARKS.md} (50 M codes, 32
threads, AVX2-only host, $K = 100$):

\begin{center}
\begin{tabular}{lrr}
\toprule
\textbf{Kernel} & \textbf{cmp/s} & \textbf{GB/s} \\
\midrule
\texttt{bench\_hamming\_top1} (NQ=8 batched) & \textbf{2308 M} & 27.7 \\
\texttt{bench\_hamming} (top-K=100) & \textbf{740 M} & 71.1 \\
\texttt{bench\_asym\_b1} (single-threaded) & 0.55 M & 0.1 \\
\texttt{bench\_asym\_b2} (single-threaded) & 2.19 M & 0.4 \\
\texttt{bench\_asym\_b4} (single-threaded) & 2.36 M & 0.9 \\
\bottomrule
\end{tabular}
\end{center}

Three observations.

\paragraph{(1) The Hamming kernel hit the bandwidth wall.} 71.1 GB/s
on the single-query top-K path is within a few percent of the dev
host's sustained DRAM bandwidth. The \texttt{NQ=8} batched top-1
reaches 27.7 GB/s because each load is reused 8$\times$, so the
effective bandwidth is $27.7 \cdot 8 = 222$ GB/s of ``consumed'' data
per second --- the chip's \emph{L3} bandwidth, since the working set
spills out of L2. The implication: at this scale, the symmetric kernel
cannot get meaningfully faster without lossless code compression (out
of scope for v1) or an outer index (out of scope for v1).

\paragraph{(2) Asymmetric $b \in \{2, 4\}$ are within striking
distance of the threaded ceiling.} The ratio $2.19$ M cmp/s $\times$
32 threads $\approx 70$ M cmp/s (assuming linear scaling, modulo
memory contention) lands at roughly \textbf{5\% of the bandwidth
headroom on this host}. That headroom collapses on AVX-512F because
the per-cycle FMA work doubles; we expect the threaded Plan-2 number
on an AVX-512F box to be \textbf{bandwidth-bound, not compute-bound},
at roughly the 200--300 M cmp/s scale that makes a 5 M corpus
searchable in $\sim 20$ ms.

\paragraph{(3) The $b = 1$ asymmetric path is the canary.} It is
reading the \emph{least} memory of any asymmetric kernel and is the
\emph{slowest}. The diagnosis (\S\ref{sec:b1-outlier}) is unambiguous
--- the AVX2 unpack --- and the fix is queued. Plan 1's contribution
is that this gap is now visible; under the previous setup it would
have looked like ``asymmetric is just slow.''

\section{The science questions Plan 1 sets up for Plan 3}

Plan 1 ships kernels; Plan 3 ships a recipe sweep over BEIR. The
specific scientific questions Plan 1's machinery is engineered to
answer:

\begin{center}
\begin{tabular}{p{0.55\textwidth}l}
\toprule
\textbf{Question} & \textbf{Recipes that answer it} \\
\midrule
Does asymmetric $b = 2$ at full $d$ match $b = 1$ + float rescore at the same byte budget? & R0 vs R3 \\
Does halving $d$ and doubling $b$ win under random-orthogonal $W$? & R3 vs R4 \\
Is $b = 4$ at $d/2$ (same byte budget as $b = 2$ at $d$) a Pareto move? & R3 vs R5 \\
How much of the float-baseline NDCG does any of the above sacrifice? & R6 (existing engines) \\
\bottomrule
\end{tabular}
\end{center}

The headline finding the team will read off the resulting CSV is
whether \textbf{asymmetric multi-bit replaces the binary-plus-rescore
pattern} that current engines rely on. If R3 wins R0 at the same byte
budget, the deployment recommendation across the embedding fleet
flips. If R0 still wins, ASH-style learned projection (v2) becomes
the next research lever; if it loses by a meaningful margin, v2 is
unnecessary.

The measurement infrastructure for that decision is already in this
tree. Plan 1 is the foundation; Plan 2 makes it callable from Python;
Plan 3 turns the dial.

\section{Limitations of Plan 1 (read before citing)}

\begin{enumerate}
    \item \textbf{AVX-512F throughput is unmeasured.} The dev host
    lacks AVX-512F; correctness is tested via the \texttt{*\_avx2}
    ctest variant that forces the fallback path on the same source.
    Numbers in \S\ref{sec:results} are AVX2-only.
    \item \textbf{Asymmetric scans are single-threaded.} The OpenMP
    shard/merge pattern from \texttt{scan\_batch\_parallel\_topk} will
    be reused in Plan 2. Multiply the asymmetric $b = 2$/$b = 4$
    numbers by $\sim 16$--$24\times$ to approximate the threaded
    ceiling at 32 threads (modulo memory bandwidth, which on AVX2-only
    hosts becomes the binding constraint above $\sim 16$ threads).
    \item \textbf{The AVX2 $b = 1$ unpack is suboptimal}
    (\S\ref{sec:b1-outlier}). The headline $b = 1$ cmp/s number
    under-represents what the kernel will look like once the
    \texttt{vpshufb}-based unpack lands.
    \item \textbf{No $D \ne 768$.} Plan 1 hard-codes
    $\texttt{ASYM\_D} = 768$. Templates for $D \in \{384, 1024,
    1536\}$ come in Plan 3.
    \item \textbf{No Python, no engine.} \texttt{SimdqIndex} and
    \texttt{SimdqEngine} are Plan 2 / Plan 3. Today the kernels are
    callable only from C.
\end{enumerate}

\section*{References}

\begin{itemize}[leftmargin=*]
    \item Charikar, M. (2002). ``Similarity Estimation Techniques from
    Rounding Algorithms.'' STOC. --- The SimHash bound (\S2); states
    Eq.~(1) as a corollary of Goemans--Williamson rounding.
    \item Goemans, M.\ X.\ \& Williamson, D.\ P.\ (1995). ``Improved
    Approximation Algorithms for Maximum Cut and Satisfiability
    Problems Using Semidefinite Programming.'' JACM 42(6),
    pp.~1115--1145. --- The original random-hyperplane rounding
    argument; Lemma~3.2 contains the geometric identity used in
    Appendix~\ref{app:simhash-proof}.
    \item Vempala, S.\ (2004). \emph{The Random Projection Method}.
    AMS DIMACS Series, vol.~65. --- Textbook treatment of random
    projections, including the SimHash identity.
    \item Andoni, A.\ \& Indyk, P.\ (2008). ``Near-Optimal Hashing
    Algorithms for Approximate Nearest Neighbor in High Dimensions.''
    Communications of the ACM 51(1), pp.~117--122. --- LSH context
    and survey.
    \item Tepper, M.\ \& Willke, T. (2026, internal pre-print).
    ``Asymmetric Scalar Hashing.'' --- The asymmetric $b$-bit +
    per-vector scale framework Plan 1 ports (\S4).
    \item Mu\l{}a, W., Kurz, N., \& Lemire, D. (2018). ``Faster
    Population Counts using AVX2 Instructions.'' Computer Journal. ---
    The nibble-LUT popcount used in
    \path{simdq_kernels_hamming.h:135-140}.
    \item Existing \texttt{embedding\_hashes/hb5} driver in
    \path{~/sandbox2}. --- Source of the NQ-batched query trick
    (\S3.4).
\end{itemize}

\appendix

\section{Proof of the Charikar SimHash bound}\label{app:simhash-proof}

We prove the per-hyperplane identity stated as Eq.~(1) in \S2:
\begin{equation}\label{eq:simhash-restate}
\Pr_{h \sim \mathcal{N}(0, I_D)}\bigl[\,\mathrm{sign}(\langle h, x\rangle) \neq \mathrm{sign}(\langle h, y\rangle)\,\bigr]
\;=\; \frac{\theta(x, y)}{\pi},
\end{equation}
for unit vectors $x, y \in \mathbb{R}^D$ with angle $\theta = \theta(x, y) = \arccos\langle x, y\rangle \in [0, \pi]$. The agreement form in Eq.~(1) follows by complementation.

The argument is the Goemans--Williamson random-hyperplane rounding identity (Goemans \& Williamson 1995, Lemma~3.2), applied to a pair of unit vectors. We give it in three short steps.

\paragraph{Step 1: reduce to the 2D plane spanned by $x$ and $y$.} Let $P \subset \mathbb{R}^D$ be the 2-dimensional subspace spanned by $x$ and $y$ (assume $\theta \in (0, \pi)$, i.e.\ $x, y$ are linearly independent; the boundary cases $\theta \in \{0, \pi\}$ give probability $0$ and $1$ respectively, consistent with~\eqref{eq:simhash-restate}). Decompose
\[
h \;=\; h_P + h_{P^\perp},
\]
where $h_P$ is the orthogonal projection of $h$ onto $P$ and $h_{P^\perp} \perp P$. Since $x, y \in P$, we have $\langle h, x\rangle = \langle h_P, x\rangle$ and $\langle h, y\rangle = \langle h_P, y\rangle$, so the event in~\eqref{eq:simhash-restate} depends only on $h_P$.

By rotational invariance of the standard Gaussian on $\mathbb{R}^D$, $h_P$ is itself an isotropic Gaussian on the 2D subspace $P$, i.e.\ $h_P \sim \mathcal{N}(0, I_2)$ in any orthonormal basis of $P$. The problem therefore reduces to the 2D case.

\paragraph{Step 2: in 2D, the direction of $h_P$ is uniform on $S^1$.} Write $h_P = R \cdot u$, where $R = \|h_P\| \geq 0$ and $u = h_P / \|h_P\| \in S^1$ (defined almost surely, since $h_P = 0$ has probability zero). For an isotropic 2D Gaussian, $R$ and $u$ are independent and $u$ is \textbf{uniform on the unit circle} $S^1$ --- this is the standard polar decomposition of $\mathcal{N}(0, I_2)$, equivalent to the Box--Muller construction.

The signs $\mathrm{sign}(\langle h_P, x\rangle) = \mathrm{sign}(R \langle u, x\rangle) = \mathrm{sign}(\langle u, x\rangle)$ depend only on $u$, not on $R$. So
\[
\Pr_h\bigl[\mathrm{sign}(\langle h, x\rangle) \neq \mathrm{sign}(\langle h, y\rangle)\bigr]
\;=\;
\Pr_{u \sim \mathrm{Unif}(S^1)}\bigl[\mathrm{sign}(\langle u, x\rangle) \neq \mathrm{sign}(\langle u, y\rangle)\bigr].
\]

\paragraph{Step 3: count the separating arc.} Choose an orthonormal basis of $P$ in which $x = (1, 0)$ and $y = (\cos\theta, \sin\theta)$. Parametrise $u \in S^1$ by its angle $\phi \in [0, 2\pi)$, so $u = (\cos\phi, \sin\phi)$. Then
\[
\langle u, x\rangle = \cos\phi, \qquad \langle u, y\rangle = \cos(\phi - \theta).
\]
Each inner product changes sign exactly twice as $\phi$ traverses $[0, 2\pi)$:
\begin{itemize}
  \item $\cos\phi$ flips sign at $\phi = \pi/2$ and $\phi = 3\pi/2$;
  \item $\cos(\phi - \theta)$ flips sign at $\phi = \theta + \pi/2$ and $\phi = \theta + 3\pi/2$.
\end{itemize}
The four sign-change points partition $S^1$ into four arcs. Walking around the circle, the two signs disagree on exactly the two arcs of length $\theta$ each (one between $\pi/2$ and $\theta + \pi/2$, one between $3\pi/2$ and $\theta + 3\pi/2$, modulo $2\pi$). The total angular measure of disagreement is therefore $2\theta$, out of total measure $2\pi$. Since $u$ is uniform:
\[
\Pr_u\bigl[\mathrm{sign}(\langle u, x\rangle) \neq \mathrm{sign}(\langle u, y\rangle)\bigr]
\;=\; \frac{2\theta}{2\pi} \;=\; \frac{\theta}{\pi},
\]
which is~\eqref{eq:simhash-restate}. \hfill$\blacksquare$

\paragraph{Remarks.}
\begin{itemize}
  \item The proof never used $D$. The bound is dimension-free: the probability depends only on the angle $\theta$, regardless of the ambient dimension.
  \item The same argument shows that the LHS depends on $x$ and $y$ only through $\langle x, y\rangle$ (since rotations of $\mathbb{R}^D$ leave both the Gaussian and the inner product invariant). This is what makes Hamming distance on the binary codes a function of cosine alone.
  \item The bound is \emph{tight} --- it is an exact identity, not an inequality. The looseness in practice comes from the variance of the empirical bit-agreement rate around its mean (Eq.~(2) in \S2), not from the bound itself.
  \item For $\theta$ small (the regime of interest in retrieval), $\theta/\pi \approx (1 - \cos\theta)^{1/2} \cdot \sqrt{2}/\pi$, so the disagreement probability scales as $\sqrt{1 - \cos\theta}$ near $\theta = 0$ --- bit codes resolve near-duplicate similarities sub-linearly in $1 - \cos\theta$, which is the underlying reason the Hamming top-1 estimator can rank tightly clustered candidates.
\end{itemize}

\end{document}
